% ============================================================
% Sección 4: Resultados
% ============================================================

\section{Resultados}
\label{sec:resultados}

\subsection{Validación 1D - Helmholtz con solución analítica}
\label{subsec:resultados1d}

La Tabla~\ref{tab:resultados1d} muestra las métricas de precisión obtenidas para la
ecuación de Helmholtz 1D con la arquitectura SIREN $5\times64$, $\omega_0 = 1.0$.

\begin{table}[t]
\centering
\renewcommand{\arraystretch}{1.2}
\caption{métricas de precisión, PINN-SIREN 1D ($k = 2\pi$)}
\label{tab:resultados1d}
\begin{tabular}{lc}
\hline
Métrica & Valor \\
\hline
Error L2 (\%)          & 0.006 \\
MSE                    & $1.5 \times 10^{-9}$ \\
RMSE                   & $3.9 \times 10^{-5}$ \\
MAE                    & $3.50 \times 10^{-5}$ \\
Error máximo           & $5.77 \times 10^{-5}$ \\
$R^2$                  & 1.000000 \\
Correlación de Pearson & 1.000000 \\
\hline
Épocas Adam            & 15{,}000 / 15{,}000 \\
Evaluaciones L-BFGS    & 202 / 500 \\
Tiempo total           & 131.1 s \\
\hline
\end{tabular}
\end{table}
\noindent\footnotesize Evaluado en malla uniforme de 1{,}000 puntos en $[0,1]$.
Arquitectura SIREN: 5 capas $\times$ 64 neuronas, $\omega_0 = 1.0$.
Solución analítica: $E_{\mathrm{exacta}}(x) = \cos(kx)$.
GPU: NVIDIA RTX 5050, SEED = 42, laptop conectada a corriente y sin
procesos concurrentes. ``Evaluaciones L-BFGS'' cuenta las llamadas a la
función de pérdida durante la búsqueda de línea \textit{strong Wolfe}, no
iteraciones del optimizador. En una corrida previa sin estas condiciones el
tiempo total fue 504.2 s: el mismo hardware, en régimen térmico degradado,
reprodujo idénticamente todas las métricas de precisión a 60.4 ms/época
frente a 8.75 ms/época en condiciones controladas.
\normalsize

El error L2 de $0.006\%$ se mantiene por debajo del umbral de la tesis ($< 5\%$) por
casi tres órdenes de magnitud.
El ajuste es prácticamente perfecto: $R^2 = 1.000000$ y correlación de Pearson $= 1.000000$
en toda la malla de evaluación.

\begin{figure}[htbp]
  \centering
  \includegraphics[width=0.75\linewidth]{resultados_pinn_1d.png}
  \caption{Validación PINN-SIREN para Helmholtz 1D.
           Panel A: comparación entre la predicción de la red y la solución
           analítica $E_{\mathrm{exacta}}(x) = \cos(kx)$ en 1{,}000 puntos
           uniformes en $[0,1]$. Panel B: error puntual
           $|\hat{E}(x) - E_{\mathrm{exacta}}(x)|$. Panel C: curva de
           pérdida total durante las fases Adam y L-BFGS.
           Error $L^2 = 0.006\%$, $R^2 = 1.000000$.
           GPU: NVIDIA RTX 5050, SEED = 42.}
  \label{fig:resultados1d}
\end{figure}

\begin{figure}[htbp]
  \centering
  \includegraphics[width=0.9\linewidth]{metricas_adicionales_1d.png}
  \caption{Métricas adicionales de entrenamiento NB01. Panel A:
           distribución del error absoluto
           $|\hat{E}(x) - E_{\mathrm{exacta}}(x)|$. Panel B: error relativo
           puntual (\%). Panel C: dispersión de la predicción PINN contra
           la solución exacta, con la línea $y = x$ como referencia de
           ajuste perfecto.}
  \label{fig:metricas1d}
\end{figure}

\subsection{Validación 2D - Helmholtz con campo complejo}
\label{subsec:resultados2d}

La Tabla~\ref{tab:resultados2d} reporta las métricas para la validación 2D con la
arquitectura optimizada SIREN $5\times128$, $\omega_0 = 1.0$, y muestreo LHS de
3{,}000 puntos.

\begin{table}[t]
\centering
\footnotesize
\renewcommand{\arraystretch}{1.2}
\setlength{\tabcolsep}{4pt}
\caption{métricas de precisión, PINN-SIREN 2D ($k = 2\pi$, $k_x = k_y = k/\sqrt{2}$)}
\label{tab:resultados2d}
\begin{tabular}{@{}lcc@{}}
\hline
Métrica & $E_{\mathrm{real}}$ & $E_{\mathrm{imag}}$ \\
\hline
Error L2 (\%)          & 0.214 & 0.127 \\
MSE                    & $2.20 \times 10^{-6}$ & $8.43 \times 10^{-7}$ \\
RMSE                   & $1.48 \times 10^{-3}$ & $9.18 \times 10^{-4}$ \\
MAE                    & $1.27 \times 10^{-3}$ & $7.71 \times 10^{-4}$ \\
Error máximo           & $3.05 \times 10^{-3}$ & $2.93 \times 10^{-3}$ \\
$R^2$                  & 0.999995 & 0.999998 \\
Correlación de Pearson & 0.999998 & 0.999999 \\
\hline
Error L2 promedio (\%) & \multicolumn{2}{c}{0.171} \\
\hline
\end{tabular}
\end{table}
\noindent\footnotesize Evaluado en malla uniforme de $100 \times 100$ puntos.
Arquitectura: SIREN $5 \times 128$, $\omega_0 = 1.0$, 66{,}690 parámetros.
Puntos de colocación: LHS $N_c = 3{,}000$.
Frontera: onda plana $e^{i(k_x x + k_y y)}$ en 4 bordes, $N_b = 300$ pts/borde.
Entrenamiento: Adam (8{,}737 épocas) + L-BFGS (1{,}035 evaluaciones de la
función objetivo), tiempo total 209.9 s (GPU laptop conectada a corriente,
sin procesos concurrentes).
Dispersión entre semillas (SEED $\in \{42, 123, 777\}$), medidas en la
misma sesión:
error $L^2$ promedio $= 0.192 \pm 0.089$\,\% (media $\pm$ desv.\ est.).
\normalsize

El error L2 promedio de $0.171\%$ se mantiene por debajo del umbral de la tesis
($< 5\%$) por más de un orden de magnitud. La precisión se mantiene alta en ambas
componentes del campo complejo, con $R^2 > 0.9999$ para $E_{\mathrm{real}}$ y
$E_{\mathrm{imag}}$.

En un experimento preliminar con $d = 64$ (mismo $N_c$, $N_b$ y protocolo de
optimización) el error promedio fue $0.436\%$, frente a $0.171\%$ con
$d = 128$. Esta comparación involucra además una diferencia en el número de
evaluaciones de L-BFGS entre ambas corridas, por lo que la mejora no puede
atribuirse de forma aislada al aumento de capacidad de la red; se reporta
como evidencia cualitativa, no como una ablación controlada. No se dispone
de un experimento que aísle la contribución del muestreo LHS frente a una
malla cartesiana equivalente.

\begin{figure}[htbp]
  \centering
  \includegraphics[width=0.75\linewidth]{resultados_pinn_2d.png}
  \caption{Validación PINN-SIREN para Helmholtz 2D con campo complejo.
           Fila superior (A1-A3): solución analítica, predicción PINN-SIREN
           y error absoluto para $E_{\mathrm{real}}$. Fila central (B1-B3):
           lo mismo para $E_{\mathrm{imag}}$. Fila inferior: intensidad
           predicha $I = |E|^2$ (C1), puntos de colocación LHS (C2) y
           curvas de pérdida durante Adam y L-BFGS (C3).
           La diferencia visual entre solución exacta y predicción es
           indistinguible; los errores puntuales son del orden
           de $10^{-3}$ (ver Tabla~\ref{tab:resultados2d}).
           Malla de evaluación $100 \times 100$, SEED = 42.}
  \label{fig:resultados2d}
\end{figure}

\begin{figure}[htbp]
  \centering
  \includegraphics[width=0.9\linewidth]{metricas_adicionales_2d.png}
  \caption{Métricas adicionales de entrenamiento NB02. Paneles A-B:
           distribución del error absoluto
           $|E_{\mathrm{real}} - \hat{E}_{\mathrm{real}}|$ y
           $|E_{\mathrm{imag}} - \hat{E}_{\mathrm{imag}}|$. Paneles C, F:
           dispersión de la predicción PINN contra la solución exacta para
           cada componente. Paneles D-E: error relativo puntual (\%) en la
           malla $100 \times 100$; los valores más altos se concentran en
           bandas antidiagonales ($x+y\approx$ const.) donde la solución
           exacta cruza por cero, no en las esquinas del dominio. La curva
           de pérdida de entrenamiento se muestra en el Panel C3 de la
           Figura~\ref{fig:resultados2d}.}
  \label{fig:metricas2d}
\end{figure}

\subsection{Contexto en la literatura}
\label{subsec:comparativa}

La Tabla~\ref{tab:comparativa} sitúa los resultados de precisión de la PINN-SIREN
propuesta junto al antecedente más cercano en la literatura.

\begin{table*}[t]
\centering
\renewcommand{\arraystretch}{1.2}
\caption{error L2 relativo, PINN-SIREN y antecedente de la literatura}
\label{tab:comparativa}
\begin{tabular}{lccc}
\hline
Método & Dimensión & Comparación & Error L2 (\%) \\
\hline
Schoder y Kraxberger \cite{Schoder2024_PINN_acoustics} & 3D (acústica) & PINN vs.\ FEM (malla fina) & 2.49 \\
Zhang et al. \cite{Zhang2025_FEPIRBN} & 2D (EM) & PINN vs.\ simulación numérica & 1.40--5.82 \\
\textbf{PINN-SIREN (este trabajo)} & \textbf{1D (óptica)} & PINN vs.\ solución analítica & \textbf{0.006} \\
\textbf{PINN-SIREN (este trabajo)} & \textbf{2D (óptica)} & PINN vs.\ solución analítica & \textbf{0.171} \\
\hline
\end{tabular}

\vspace{4pt}
\begin{minipage}{\textwidth}
\footnotesize Los cuatro resultados no son directamente comparables entre sí: difieren en
dimensionalidad, dominio físico (acústico vs.\ óptico) y tipo de solución
de referencia (FEM vs.\ analítica). Se presentan como contexto de la
literatura, no como línea base de comparación cuantitativa. Los resultados
propios se validan contra el umbral fijo de la tesis ($L^2 < 5\%$),
independiente de trabajos externos.
\end{minipage}
\end{table*}
